\chapter{Algebraic code skeletons from Koszul duality}
\label{chap:holographic-codes-koszul}
\index{holographic code!algebraic skeleton from Koszul duality|textbf}
\index{bar-cobar recovery!Koszul duality|textbf}

\providecommand{\KL}{\mathrm{KL}}
\providecommand{\QEC}{\mathrm{QEC}}

Fix a chart $b$ in a factorisation dg-category
$\mathcal C^{\mathrm{op}}$ on $(X,D,\tau)$ and a connected positive-weight
quadratic presentation
$\cA=A_b=T_X(V)/(R)$.  The recovery problem involves five objects:
\[
\cA = A_b = \End_\mathcal{C}(b)\quad\text{(level 1, chart algebra),}
\qquad
B_X(\cA) = T^c(s^{-1}\bar\cA)\quad\text{(level 2, full bar coalgebra),}
\]
\[
\cA^{\mathrm i}=C_X(s^{-1}V,s^{-2}R)\subset T^c(s^{-1}V)
\quad\text{(quadratic coalgebra),}
\qquad
\cA^!_\infty=\mathbb D_{\operatorname{Ran}}B_X(\cA)
\quad\text{(post-Verdier algebra under }H_{\mathbb D}(\cA)\text{),}
\qquad
Z^{\mathrm{der}}_{\mathrm{ch}}(\cA) = C^\bullet_{\mathrm{ch}}(\cA,\cA)
\quad\text{(level 3, universal closed sector).}
\]
Local finite Verdier duality also gives the strict quadratic algebra
$\cA^!=\mathbb D_{\operatorname{Ran}}\cA^{\mathrm i}$.

Three typed comparison lanes govern these objects.  First, enhanced
bar--cobar duality in the pro-nilpotent Ran category gives the universal
counit
\[
 \varepsilon_\cA\colon\Omega_XB_X(\cA)\xrightarrow{\sim}\cA
\]
for every augmented algebra in the ambient of Theorem~A; factorisation
and completed chain models carry the packages
$H_{\mathrm{fact}}$ and $H_{\mathrm{conv}}$.  This is the Ran form of
Francis--Gaitsgory~\cite[Proposition~4.1.2]{FrancisGaitsgory2012} and
the ordinary resolution of
Loday--Vallette~\cite[Corollary~2.3.4]{LodayVallette12}.
Second, the chosen presentation determines the canonical comparison
and its adjoint
\[
 q_\cA\colon\cA^{\mathrm i}\longrightarrow B_X(\cA),
 \qquad
 p_\cA\colon\Omega_X(\cA^{\mathrm i})\longrightarrow\cA.
\]
Under $H_{\mathrm{CL}}(\cA,\cA^{\mathrm i},\tau_{\mathrm i})$,
Theorem~B identifies chiral Koszulness with either map being a
quasi-isomorphism; this is the quadratic criterion modeled by
Loday--Vallette~\cite[Theorem~3.4.6]{LodayVallette12}.  Third, for a
fixed coalgebra $C$, Positselski's equivalence has the typed form
\[
 \mathsf D^{\mathrm{co}}(C\text{-}\mathsf{Comod})
 \simeq
 \mathsf D^{\mathrm{ctr}}(C\text{-}\mathsf{Contramod}).
\]
This is Positselski~\cite[Theorem~5.2]{Positselski11}.
Its carrier is the module theory of one fixed $C$; the quadratic
recognition carrier is the coalgebra map $q_\cA$.

The level-$3$ object is the derived chiral centre
$Z^{\mathrm{der}}_{\mathrm{ch}}(\cA)$, equipped with its universal
open/closed action by Chapter~\ref{ch:holographic-datum-master} and
Theorem~\ref{thm:thqg-swiss-cheese}.  A physical realization datum for
a chosen holomorphic-topological theory~$T$ is a comparison
\[
 \beta_T\colon
 \operatorname{Loc}_{\partial}\operatorname{Obs}^{\mathrm{bulk}}(T)
 \longrightarrow Z^{\mathrm{der}}_{\mathrm{ch}}(\cA).
\]
Thus $\varepsilon_\cA$, $q_\cA$, $p_\cA$, and $\beta_T$ have distinct
sources, targets, and hypothesis packages; $H_{\mathbb D}(\cA)$ adds
the Verdier transpose $\mathbb D_{\operatorname{Ran}}(q_\cA)$.
Chapter~\ref{chap:entanglement-modular-koszul} reads the
Calabrese--Cardy entanglement entropy (proved elsewhere) through the
scalar projection of this structure; every reconstruction beyond that
scalar shadow is conditional on its named comparison data.
The algebraic recovery statement is:

\begin{center}
\fbox{\parbox{0.85\textwidth}{\centering
\textbf{Chiral Koszulness is quadratic exactness: the comparison
$q_\cA\colon\cA^{\mathrm i}\to B_X(\cA)$, equivalently
$p_\cA\colon\Omega_X(\cA^{\mathrm i})\to\cA$, is a
quasi-isomorphism.}  The universal counit $\varepsilon_\cA$ resolves
every augmented chart in the Theorem~A ambient.  Theorem~C supplies a
Verdier--Lagrangian algebraic code skeleton.  Unitary completion and a
physical error algebra supply Hilbert-space quantum error correction;
the comparison $\beta_T$ supplies physical bulk realization.
}}
\end{center}

\noindent
The characterizations of Koszulness (K1--K12,
Theorem~\ref{thm:koszul-equivalences-meta}: seven intrinsic
bidirectional equivalences, one listed consequence
(Barr--Beck--Lurie monadicity), one conditional
factorization-homology comparison, one Theorem~H-conditional chiral
Hochschild consequence, one conditional on the shifted-symplectic
hypothesis package, and one one-directional
$\mathcal D$-module purity implication) yield corresponding diagnostics
for the quadratic comparison $q_\cA$.
The shadow depth classifies the non-scalar layers in the formal
MC obstruction tower.
The Lagrangian isotropy of Theorem~C provides the code/complement
decomposition: each summand is isotropic and the cross-pairing is
perfect.  The derived-centre/BRST comparison and the corresponding
Hilbert-space inner product promote this algebraic polarization to a
physical holographic recovery problem.

% ======================================================================
%
% 1. SYMPLECTIC CODE STRUCTURE FROM LAGRANGIAN ISOTROPY
%
% ======================================================================

\section{Symplectic code structure from Lagrangian isotropy}
\label{sec:hc-knill-laflamme}
\index{symplectic code!from Lagrangian isotropy|textbf}

\begin{remark}[Three algebraic recovery structures]
\label{rem:hc-two-codes}
\index{holographic code!three recovery structures}
The algebraic framework produces three related recovery structures:
\begin{enumerate}[label=\textup{(\alph*)}]
\item The \emph{linearized symplectic code}: $\cC = \mathbf{Q}_g(\cA) \subset
 \mathbf{C}_g(\cA) = \cH$
 (Lagrangian summand inside the ambient complex).
 This is the algebraic structure relevant to a Knill--Laflamme
 comparison after a physical inner product is supplied.
\item The \emph{universal resolution}: the bar functor sends $\cA$
 to $B_X(\cA)$, and the counit
 $\varepsilon_\cA\colon\Omega_XB_X(\cA)\to\cA$ is an equivalence in
 the Theorem~A ambient for every augmented chart.
\item The \emph{quadratic compression}: the presentation
 $T_X(V)/(R)$ replaces the full bar coalgebra by
 $\cA^{\mathrm i}=C_X(s^{-1}V,s^{-2}R)$.  Its decoding map is
 $p_\cA\colon\Omega_X(\cA^{\mathrm i})\to\cA$, and Theorem~B makes
 its exactness equivalent to $q_\cA$ being a quasi-isomorphism.
\end{enumerate}
The linearized symplectic code governs the genus-by-genus phase-space
polarization.  The universal resolution reconstructs the full chart.
The quadratic compression measures the efficiency of the chosen
presentation.  The Lagrangian splitting in~(a) is the conditional
Theorem~C upgrade with its own perfectness and non-degeneracy package.
\end{remark}

\begin{proposition}[Verdier isotropy and the Knill--Laflamme obstruction;
\ClaimStatusConditional]
\label{prop:hc-knill-laflamme}
\index{Knill--Laflamme condition|textbf}
\index{Lagrangian isotropy!Knill--Laflamme}
\textit{Type signature.} Quadrant: Open. Presentation: chain-level.
Level: $5$ (scalar genus-$g$ pairing on the conformal-block carrier).
Hypothesis package: chirally Koszul pair $(\cA, \cA^!)$;
finite-dimensional shadow window; Verdier involution $\sigma$ as in
Proposition~\textup{\ref{prop:lagrangian-eigenspaces}}; complementary
Lagrangian polarisation of Theorem~C.
\smallskip\noindent
Let $(\cA, \cA^!)$ be a chirally Koszul pair satisfying the
hypotheses of the Lagrangian clause of
Theorem~\textup{\ref{thm:quantum-complementarity-main}} at genus
$g\geq1$, with the Verdier/Shapovalov comparison of
Proposition~\textup{\ref{prop:lagrangian-eigenspaces}}.
Define:
\begin{itemize}
\item $\cH := \mathbf{C}_g(\cA)$ \quad
 \textup{(}ambient complex\textup{)},
\item $\cC := \mathbf{Q}_g(\cA)$ \quad
 \textup{(}Lagrangian summand, the linearized code subspace\textup{)},
\item $\mathcal E := \mathbf{Q}_g(\cA^!)$ \quad
 \textup{(}complementary Lagrangian; a physical Hermitian form
 determines the associated orthogonal complement\textup{)}.
\end{itemize}
Conditional on the Lagrangian-eigenspace package in
Proposition~\textup{\ref{prop:lagrangian-eigenspaces}}, the following
two-layer symplectic structure holds:

\begin{enumerate}[label=\textup{(\roman*)}]
\item \emph{Verdier isotropy \textup{(}summand decoupling\textup{)}.}
 The Verdier pairing $\langle \cdot, \cdot \rangle_{\mathbb{D}}$
 vanishes on $\cC \times \cC$ and on
 $\mathcal E \times \mathcal E$:
 \begin{equation}\label{eq:hc-verdier-isotropy}
 \langle v, w \rangle_{\mathbb{D}} = 0
 \quad\text{for } v, w \in \cC
 \;\;\text{(resp.\ both in } \mathcal E\text{)}.
 \end{equation}
 This is the algebraic null-pairing carried by a symplectic code
 model.  The Hilbert-space Knill--Laflamme equation uses, in addition,
 a Hermitian adjoint and a specified physical error algebra.

\item \emph{Perfect Shapovalov cross-pairing.}
 Define $\langle v, w \rangle_S := \langle v, \sigma(w)
 \rangle_{\mathbb{D}}$.
 Since $\cC = V^+$ and $\mathcal E = V^-$ are the $\pm1$-eigenspaces
 of the Verdier involution~$\sigma$
 \textup{(}Proposition~\textup{\ref{prop:lagrangian-eigenspaces}(iii)}\textup{)},
 the Shapovalov form restricts to:
 \begin{equation}\label{eq:hc-shapovalov-cross-pairing}
 \langle v, w \rangle_S = -\langle v, w \rangle_{\mathbb{D}}
 \quad\text{for } v \in \cC,\; w \in \mathcal E.
 \end{equation}
 The Verdier cross-pairing
 $\cC \times \mathcal E \to \bC$ is non-degenerate
 \textup{(}it is the Lagrangian polarization\textup{)},
 hence the Shapovalov cross-pairing is equally non-degenerate
 with a sign flip.
 The bipartite decomposition
 $\cH = \cC \oplus \mathcal E$ is a \emph{Lagrangian}
 splitting: its summands are Verdier-isotropic and pair perfectly
 across the decomposition.
\end{enumerate}

\smallskip\noindent
Part~(i) establishes the \emph{Lagrangian prerequisite}
for the symplectic code: the Verdier form vanishes internally on each
summand.
Part~(ii) shows the cross-pairing is non-degenerate
with a sign flip, giving a \emph{symplectic}
code/complement structure.
The full Knill--Laflamme condition
$P_\cC E_i^\dagger E_j P_\cC = \lambda_{ij} P_\cC$
for a physical error family $\{E_i\}$
requires, in addition, the physical Hilbert space inner product
\textup{(}from the state--operator correspondence
in unitary theories\textup{)}, which is a separate
structure from the Verdier pairing.
On a one-dimensional scalar conformal-block line the KL compression
condition is automatic.  The full genus-$1$ complementarity complex
retains the dimension supplied by its conformal-block carrier.
\end{proposition}

\begin{proof}
\emph{Part (i).}
The Verdier involution $\sigma$ satisfies
$\langle \sigma(v), \sigma(w) \rangle_{\mathbb{D}} =
-\langle v, w \rangle_{\mathbb{D}}$
(anti-commutativity,
Proposition~\ref{prop:lagrangian-eigenspaces}(ii)).
For $v, w \in \cC = V^+$ (eigenvalue $+1$ of $\sigma$):
$\langle v, w \rangle_{\mathbb{D}} =
\langle \sigma(v), \sigma(w) \rangle_{\mathbb{D}} =
-\langle v, w \rangle_{\mathbb{D}}$,
so $\langle v, w \rangle_{\mathbb{D}} = 0$.
The same argument with eigenvalue $-1$ gives
isotropy of $\mathcal E$.

\emph{Part (ii).}
For $v \in \cC = V^+$ and $w \in \mathcal E = V^-$:
$\sigma(v) = v$ and $\sigma(w) = -w$, so
$\langle v, w \rangle_S = \langle v, \sigma(w)
\rangle_{\mathbb{D}} = \langle v, -w \rangle_{\mathbb{D}}
= -\langle v, w \rangle_{\mathbb{D}}$.
Since $V^+$ and $V^-$ are complementary Lagrangians,
the Verdier cross-pairing
$\langle -, - \rangle_{\mathbb{D}}\colon V^+ \times V^- \to \bC$
is a perfect pairing
\textup{(}non-degeneracy of the total form plus isotropy
of each factor\textup{)}.
Hence $\langle v, w \rangle_S = -\langle v, w \rangle_{\mathbb{D}}$
is equally non-degenerate.
\end{proof}

\begin{remark}[Relation to Almheiri--Dong--Harlow]
\label{rem:hc-adh-comparison}
\index{Almheiri--Dong--Harlow}
The Lagrangian decomposition
$\cH = \cC \oplus \mathcal E$
has the same bipartite form as the
Almheiri--Dong--Harlow \cite{Almheiri-Dong-Harlow15}
entanglement-wedge/complementary-wedge decomposition.  At the present
level it is an algebraic model; a Hilbert-space completion and
subregion recovery maps furnish its physical realization.
The splitting is \emph{symplectic} (Verdier-isotropic
summands with non-degenerate cross-pairing,
\textup{(\ref{eq:hc-shapovalov-cross-pairing})}).
A class
$\Phi\in Z^{\mathrm{der}}_{\mathrm{ch}}(\cA) =
C^\bullet_{\mathrm{ch}}(\cA,\cA)$
in the level-$3$ derived chiral centre is assigned to the $\cC$-side
by a chosen open/closed action and a support condition for that action.
Vanishing of the induced Verdier
pairing against $\mathcal E$ is an algebraic support condition.  A
physical comparison theorem identifies this support condition with
entanglement-wedge localization.
\end{remark}

\begin{theorem}[Symplectic code structure; \ClaimStatusConditional]
\label{thm:hc-symplectic-code}
\index{symplectic code|textbf}
\index{holographic code!symplectic structure|textbf}
\textit{Type signature.} Quadrant: Open. Presentation: chain-level.
Level: $5$ (scalar carrier $\mathbf{C}_g(\cA)$ as a
genus-$g$ symplectic vector space). Hypothesis package: chirally
Koszul pair $(\cA, \cA^!)$; perfectness/non-degeneracy of the
Verdier pairing; finite-dimensional Lagrangian window of Theorem~C
(reconstruction theorem: Theorem~C, level $5$, Lagrangian clause).
\smallskip\noindent
Let $(\cA, \cA^!)$ be a chirally Koszul pair satisfying the
hypotheses of the Lagrangian clause of
Theorem~\textup{\ref{thm:quantum-complementarity-main}}.
The Lagrangian decomposition
$\mathbf{C}_g(\cA) = \mathbf{Q}_g(\cA) \oplus \mathbf{Q}_g(\cA^!)$
defines a \emph{symplectic code} in the following precise sense:
\begin{enumerate}[label=\textup{(\roman*)}]
\item \emph{Lagrangian code subspace.}
 The code subspace $\cC := \mathbf{Q}_g(\cA)$ is Lagrangian:
 it is Verdier-isotropic
 \textup{(}Proposition~\textup{\ref{prop:hc-knill-laflamme}(i)}\textup{)}
 and, in each finite-dimensional shadow window, half-dimensional:
 $\dim \cC = \tfrac{1}{2}\dim \mathbf{C}_g(\cA)$.

\item \emph{Symplectic projection.}
 The Lagrangian splitting determines
 the symplectic projection
 \[
 \pi_\cC \colon \mathbf{C}_g(\cA) \longrightarrow
 \mathbf{Q}_g(\cA)
 \]
 along $\mathbf{Q}_g(\cA^!)$.  For any algebraic operator $E$
 preserving the code summand \textup{(}i.e.,
 $E(\cC)\subset\cC$\textup{)}, one has
 $(\pi_\cC\circ E)|_\cC=E|_\cC$.  Thus projection removes the
 $\mathbf{Q}_g(\cA^!)$ component.  A class of physical errors and a
 Hermitian inner product determine whether this projection satisfies
 the Knill--Laflamme recovery equation.

\item \emph{Non-degenerate cross-pairing.}
 The Verdier pairing restricted to
 $\cC \times \mathbf{Q}_g(\cA^!)$ is a perfect pairing
 with sign flip
 \textup{(}Proposition~\textup{\ref{prop:hc-knill-laflamme}(ii)}\textup{)},
 providing the non-degenerate duality between
 the code summand and its Lagrangian complement.
\end{enumerate}
\end{theorem}

\begin{proof}
Part~(i) follows from Theorem~C:
$\mathbf{Q}_g(\cA) + \mathbf{Q}_g(\cA^!) = \mathbf{C}_g(\cA)$
with the summands complementary Lagrangians under the Verdier
pairing. Isotropy is
Proposition~\ref{prop:hc-knill-laflamme}(i);
half-dimensionality is the Lagrangian condition.

Part~(ii): given the direct sum decomposition
$\mathbf{C}_g(\cA) = \cC \oplus \mathcal E$ with
$\mathcal E = \mathbf{Q}_g(\cA^!)$, the projection
$\pi_\cC$ is well-defined. If $E(\cC) \subset \cC$,
then $E(v)\in\cC$ for $v\in\cC$, hence
$\pi_\cC(E(v))=E(v)$.  If additionally $E|_\cC=\id_\cC$,
the projection recovers the original code vector.  The physical error
model upgrades this symplectic projection statement to a QEC recovery
theorem precisely when the Knill--Laflamme compressions are scalar.

Part~(iii) is Proposition~\ref{prop:hc-knill-laflamme}(ii):
the Verdier cross-pairing is non-degenerate by the
Lagrangian polarization, with the Shapovalov form
giving $\langle v, w \rangle_S = -\langle v, w
\rangle_{\mathbb{D}}$.
\end{proof}

\begin{proposition}[Independence of symplectic projection and
Knill--Laflamme recovery;
\ClaimStatusProvedHere]
\label{prop:hc-projection-not-kl}
\index{Knill--Laflamme condition!projection obstruction}
\textit{Type signature.} Quadrant: Open. Presentation: chain-level
(symplectic vs.\ Hermitian). Level: $5$ (Hilbert-space completion of
the scalar carrier). Hypothesis package: Verdier symplectic splitting
of Theorem~C plus an external Hilbert completion
$H_{\mathrm{phys}}$ and physical error $*$-algebra $\mathfrak E$.
\smallskip\noindent
Let $H_{\mathrm{phys}}$ be a Hilbert completion of
$\mathbf C_g(\cA)$ with orthogonal projection $P_\cC$ onto the
closure of $\cC$, and let $\mathfrak E$ be a physical error
$*$-algebra on $H_{\mathrm{phys}}$.  The Verdier splitting supplies
the algebraic projection
$\pi_\cC\colon \cC\oplus\mathcal E\to\cC$.  For
$E\in\mathfrak E$ and $v\in\cC$, write
\[
E(v)=E_{\cC}(v)+E_{\mathcal E}(v),
\qquad
E_{\cC}(v)\in\cC,\quad E_{\mathcal E}(v)\in\mathcal E.
\]
Then $\pi_\cC E(v)=E_{\cC}(v)$: the projection extracts the logical
component $E_{\cC}$.  A Knill--Laflamme code for $\mathfrak E$ is the
independent Hilbert-space assertion
\[
P_\cC E_i^\dagger E_j P_\cC=\lambda_{ij}P_\cC
\qquad(E_i,E_j\in\mathfrak E).
\]
Verdier isotropy fixes the bilinear polarization; the displayed
compression condition fixes the physical error action.
\end{proposition}

\begin{proof}
The first assertion is the definition of projection along the direct
sum $\cC\oplus\mathcal E$.  It kills the second component and leaves
the first component unchanged.  The Knill--Laflamme equation is a
statement about the Hermitian adjoint in $H_{\mathrm{phys}}$ and the
compression of products $E_i^\dagger E_j$ to the code subspace.  The
Verdier form is bilinear, shifted, and anti-symplectic under~$\sigma$.
The physical Hilbert structure supplies the Hermitian adjoint and the
compression.  Their independence is visible in the finite-dimensional
model
$V=\mathbb C e_1\oplus\mathbb C e_2\oplus
\mathbb C f_1\oplus\mathbb C f_2$ with
$\cC=\langle e_1,e_2\rangle$,
$\mathcal E=\langle f_1,f_2\rangle$, Verdier pairing
$\langle e_i,f_j\rangle_{\mathbb D}=\delta_{ij}$ and zero on
each summand, and the standard Hermitian form.  The operator
$E(e_1)=e_1$, $E(e_2)=2e_2$, $E(f_i)=0$ preserves~$\cC$ and has zero
leakage.  Its compression is
$P_\cC E^\dagger E P_\cC=\operatorname{diag}(1,4)$ on~$\cC$; the two
eigenvalues distinguish it from the scalar compression required by
Knill--Laflamme.  Thus the two structures are independent.
\end{proof}

\begin{remark}[Unitary Koszul algebras and physical QEC]
\label{rem:hc-unitary-physical-qec}
\index{Knill--Laflamme condition!unitary case}
\index{Tomita--Takesaki!modular conjugation}
For \emph{unitary} Koszul algebras
\textup{(}positive-definite Shapovalov form, physical
state--operator correspondence\textup{)},
the physical inner product provides a Hermitian
structure $\langle -, - \rangle_{\mathrm{phys}}$
on $\mathbf{C}_g(\cA)$ in addition to the Verdier
pairing. The question of whether
$\mathbf{Q}_g(\cA)$ and $\mathbf{Q}_g(\cA^!)$ are
\emph{orthogonal} in this physical inner product
is additional analytic data.  The Tomita--Takesaki analogy is read
through a chosen real form: a modular conjugation is anti-linear,
whereas the Verdier involution used in Theorem~C is a complex-linear
cochain involution.  Compatibility means that the anti-linear modular
operator implements, after complexification, the polarization defined
by the complex-linear Verdier involution.
On a one-dimensional scalar conformal-block line the KL compression is
tautological; higher-dimensional conformal-block spaces require the
full error-algebra condition of Proposition~\ref{prop:hc-projection-not-kl}.
\end{remark}

\begin{conjecture}[Physical QEC at higher genus;
\ClaimStatusConjectured]
\label{conj:hc-physical-qec}
\index{Knill--Laflamme condition!higher genus conjecture}
\textit{Type signature.} Quadrant: Open. Presentation: Hilbert-space
chain-level. Level: $5$ (genus-$g$ scalar Hilbert completion).
Hypothesis package: unitary chirally Koszul pair $(\cA,\cA^!)$;
fixed Hilbert completion $H_{\mathrm{phys}}(\cA,g)$ of $\mathbf
C_g(\cA)$; physical error $*$-algebra $\mathfrak E_g$; expected
implementation by a Tomita--Takesaki modular conjugation.
\smallskip\noindent
Let $(\cA,\cA^!)$ be a unitary chirally Koszul pair and let
$g\geq1$.  Suppose a physical Hilbert completion
$H_{\mathrm{phys}}(\cA,g)$ of $\mathbf C_g(\cA)$ and a physical
error $*$-algebra $\mathfrak E_g$ are fixed.  The physical QEC
upgrade asserts:
\begin{enumerate}[label=\textup{(\roman*)}]
\item the closures of $\mathbf Q_g(\cA)$ and
 $\mathbf Q_g(\cA^!)$ are orthogonal for
 $\langle-,-\rangle_{\mathrm{phys}}$;
\item for the orthogonal projection $P_\cC$ onto
 $\overline{\mathbf Q_g(\cA)}$,
 \[
 P_\cC E_i^\dagger E_j P_\cC
 =\lambda_{ij}P_\cC
 \qquad(E_i,E_j\in\mathfrak E_g).
 \]
\end{enumerate}
The Tomita--Takesaki expectation is that a compatible modular
conjugation supplies these two conditions in the unitary
holomorphic-topological theories arising from the open/closed
realization.  This conjecture is the physical completion problem that
extends the algebraic content of Theorem~C.
\end{conjecture}

% ======================================================================
%
% 2. BAR-COBAR RECOVERY AND DERIVED-CENTRE FUNCTORIALITY
%
% ======================================================================

\section{Bar-cobar recovery and derived-centre functoriality}
\label{sec:hc-koszulness-exact-qec}
\index{bar-cobar recovery!Koszul duality|textbf}

\begin{definition}[Typed recovery data; \ClaimStatusDefinitional]
\label{def:hc-typed-recovery-data}
Let $\cA=T_X(V)/(R)$ be a complete augmented quadratically presented
chiral algebra.  Its full reconstruction coalgebra is
\[
C_\cA=
\begin{cases}
B_X(\cA),&\text{in a finite window},\\
\widehat B_X(\cA),&\text{in the strict completed lane},
\end{cases}
\]
with universal counit
\[
\varepsilon_\cA\colon
\Omega_X^{\mathrm{cont}}(C_\cA)\longrightarrow\cA.
\]
The chosen presentation supplies
\[
\cA^{\mathrm i}=C_X(s^{-1}V,s^{-2}R),\qquad
q_\cA\colon\cA^{\mathrm i}\longrightarrow C_\cA,
\qquad
p_\cA\colon\Omega_X(\cA^{\mathrm i})\longrightarrow\cA.
\]
Here $p_\cA$ is adjoint to $q_\cA$ under the twisting-morphism
bijection.  The
\emph{derived-centre transport package} $H_Z(f)$ for a weak
equivalence $f\colon A'\to A$ consists of cofibrant chiral algebra
models and a functorial derived-centre model that sends $f$ to an
equivalence
\[
 Z^{\mathrm{der}}_{\mathrm{ch}}(f)\colon
 Z^{\mathrm{der}}_{\mathrm{ch}}(A')
 \xrightarrow{\sim}Z^{\mathrm{der}}_{\mathrm{ch}}(A).
\]
For a holomorphic-topological theory~$T$ with boundary chart~$\cA$,
the physical open/closed datum is a continuous brace dg algebra map
\[
\beta_T\colon
\bigl(\operatorname{Loc}_{\partial}
\operatorname{Obs}^{\mathrm{bulk}}(T),Q_{\mathrm{BRST}}\bigr)
\longrightarrow Z^{\mathrm{der}}_{\mathrm{ch}}(\cA)
=C^\bullet_{\mathrm{ch}}(\cA,\cA).
\]
A quasi-isomorphism $\beta_T$ in the chosen analytic completion,
together with subregion maps, realizes the universal closed sector as
the boundary-local physical bulk.  For a fixed coalgebra $C$, the
separate co/contra datum has source
$\mathsf D^{\mathrm{co}}(C\text{-}\mathsf{Comod})$ and target
$\mathsf D^{\mathrm{ctr}}(C\text{-}\mathsf{Contramod})$.
\end{definition}

\begin{proposition}[Curve-level Hochschild support and physical
transport; \ClaimStatusProvedHere]
\label{prop:hc-theorem-h-scope}
\index{Theorem H!scope in holographic codes}
\index{chiral Hochschild cohomology!physical transport}
\textit{Type signature.} Quadrant: Open. Presentation: chain-level.
Level: $3$, with source
$Z^{\mathrm{der}}_{\mathrm{ch}}(\cA)=C^\bullet_{\mathrm{ch}}(\cA,\cA)$.
Hypothesis package: $H_H(\cA;S)$ for a finite
$S\subset\mathbb Z$; the class~M lane uses its
completed/pro/$J$-adic realization; physical transport uses a
quasi-isomorphism $\beta_T$.  De Sole--Kac variational Poisson
cohomology is nonvanishing in unboundedly many degrees for
Virasoro-type Poisson vertex algebras, so the finite-support
conclusion lives strictly on the stated Koszul/bounded locus, with
the bounded-to-chart comparison part of $H_H(\cA;S)$.
\smallskip\noindent
Theorem~H gives
\[
\operatorname{Supp}\ChirHoch^\bullet(\cA,\cA)
\subseteq S.
\]
If $\beta_T$ is a quasi-isomorphism, then
\[
\operatorname{Supp}H^\bullet\bigl(
\operatorname{Loc}_{\partial}\operatorname{Obs}^{\mathrm{bulk}}(T),
Q_{\mathrm{BRST}}\bigr)
\subseteq S.
\]
The first support statement belongs to curve-level chiral Hochschild
cochains; the second is its transport through the physical comparison.
Topological and categorical Hochschild theories retain their own input
categories and comparison functors.
A perfect degree-$d$ pairing is an independent hypothesis and governs
the degree-reflection statement rather than this support transport.
\end{proposition}

\begin{proof}
The $H_H(\cA;S)$ package supplies a strong deformation retract of the
chiral Hochschild complex onto a complex supported in~$S$.  A
quasi-isomorphism $\beta_T$ identifies cohomology groups degree by
degree, and therefore transports the support inclusion to its BRST
source.
\end{proof}

\begin{theorem}[Universal reconstruction, quadratic compression, and
derived-centre transport; \ClaimStatusConditional]
\label{thm:hc-koszulness-exact-qec}
\index{G12|textbf}
\textit{Type signature.} Quadrant: Open. Presentation: complete
quadratic chiral chain model. Levels: $1\leftrightarrow2$ in
clauses~\textup{(i)--(ii)}, level~$3$ in clause~\textup{(iii)}, and
physical level~$3\to4$ in clause~\textup{(iv)}. Ambient: the
pro-nilpotent Francis--Gaitsgory chiral Ran category. Hypothesis
packages: $H_{\mathrm{fact}}$, $H_{\mathrm{conv}}$, and the finite-window or
strict completed reconstruction package of
Definition~\ref{def:finite-window-essential-image-bci} for
clause~\textup{(i)};
$H_{\mathrm{CL}}(\cA,\cA^{\mathrm i},\tau_{\mathrm i})$ for
clause~\textup{(ii)}; $H_Z$ for clause~\textup{(iii)}; a
quasi-isomorphism $\beta_T$ in the chosen analytic completion for
clause~\textup{(iv)}. Reconstruction theorems invoked: Theorem~A for
the full bar coalgebra, Theorem~B for quadratic recognition,
Theorem~H for family-indexed curve-level support, and
Theorem~\ref{thm:thqg-swiss-cheese} for the universal open/closed
action.
\smallskip\noindent
Let $\cA=T_X(V)/(R)$ and use the maps of
Definition~\textup{\ref{def:hc-typed-recovery-data}}.
\begin{enumerate}[label=\textup{(\roman*)}]
\item \emph{Universal reconstruction.}
Theorem~A gives a quasi-isomorphism
\[
\varepsilon_\cA\colon
\Omega_X^{\mathrm{cont}}(C_\cA)\xrightarrow{\sim}\cA
\]
for every augmented chart carrying the stated Ran and completion
packages.

\item \emph{Quadratic recognition.}
The following conditions are equivalent:
\[
q_\cA\colon\cA^{\mathrm i}\xrightarrow{\sim}C_\cA,
\qquad
p_\cA\colon\Omega_X(\cA^{\mathrm i})\xrightarrow{\sim}\cA.
\]
They are also equivalent to acyclicity of the left and right twisted
tensor products and to strong convergence on the quadratic PBW
diagonal.  These conditions define chiral Koszulness of the chosen
presentation.  The filtered complex
$\operatorname{Cone}(q_\cA)$ is its obstruction carrier.

\item \emph{Derived-centre transport.}
The package $H_Z(\varepsilon_\cA)$ gives
\[
Z^{\mathrm{der}}_{\mathrm{ch}}
\bigl(\Omega_X^{\mathrm{cont}}(C_\cA)\bigr)
\xrightarrow{\sim}Z^{\mathrm{der}}_{\mathrm{ch}}(\cA).
\]
On the quadratic Koszul locus, $H_Z(p_\cA)$ also gives
\[
Z^{\mathrm{der}}_{\mathrm{ch}}
\bigl(\Omega_X(\cA^{\mathrm i})\bigr)
\xrightarrow{\sim}Z^{\mathrm{der}}_{\mathrm{ch}}(\cA).
\]

\item \emph{Physical realization.}
A quasi-isomorphism $\beta_T$ identifies the completed
boundary-local BRST observable complex with
$Z^{\mathrm{der}}_{\mathrm{ch}}(\cA)$.  The subregion maps and the
physical Hermitian structure then formulate holographic recovery and
the Knill--Laflamme condition.
\end{enumerate}

For a Koszul pair satisfying the respective Theorem~C, shadow, and
modular packages, the resulting algebraic data are:
\begin{enumerate}[label=\textup{(\alph*)}]
\item the Lagrangian decomposition
$\cH_g=\cC\oplus\mathcal E$ with perfect Verdier cross-pairing;
\item the formal shadow slots classified by $r_{\max}(\cA)$;
\item the scalar information functional of
Theorem~\ref{thm:ent-scalar-entropy}.
\end{enumerate}
The physical error algebra converts the first datum into KL
compressions, assigns operational meaning to the shadow slots, and
determines Hilbert-space code distance.
\end{theorem}

\begin{proof}
Clause~\textup{(i)} is the enhanced Ran bar--cobar equivalence of
Francis--Gaitsgory and its finite-window realization, modeled over a
point by Loday--Vallette~\cite[Corollary~2.3.4]{LodayVallette12}.
Clause~\textup{(ii)} is
Theorem~\ref{thm:bar-cobar-inversion-qi}, modeled by the fundamental
twisting theorem and its quadratic specialization
\cite[Theorems~2.3.2 and~3.4.6]{LodayVallette12}.  The comparison
lemma identifies the mapping cone and the two twisted tensor
complexes as equivalent acyclicity tests.  Clause~\textup{(iii)} is
the homotopy invariance stipulated by $H_Z$.  Clause~\textup{(iv)} is
the defining consequence of the supplied physical comparison.
The final three data are the conclusions of Theorem~C,
Theorem~\ref{thm:hc-shadow-redundancy}, and
Theorem~\ref{thm:ent-scalar-entropy}, respectively.
\end{proof}

\begin{remark}[The quadratic obstruction carrier]
\label{rem:hc-non-koszul-failure}
\index{quadratic Koszulness!obstruction carrier}
Under $H_{\mathrm{CL}}$, the chosen presentation is Koszul precisely
when the filtered complexes
\[
\operatorname{Cone}(q_\cA),\qquad
\operatorname{Cone}(p_\cA)
\]
are acyclic.  Admissible and critical quotients
$L_k(\mathfrak g)$ are decided by computing these cones, or
equivalently the off-diagonal PBW groups, in the relevant completion.
The universal counit $\varepsilon_\cA$ remains the Theorem~A
resolution throughout this test.  A physical information-loss claim
is carried by the further map $\beta_T$ and its subregion recovery
data.
\end{remark}

% ======================================================================
%
% 3. SHADOW DEPTH AND FORMAL SLOTS
%
% ======================================================================

\section{Shadow depth as formal slot structure}
\label{sec:hc-redundancy}
\index{shadow depth!formal slot structure|textbf}
\index{shadow slots|textbf}

\begin{conjecture}[Shadow depth and filtered recovery models]
\ClaimStatusConjectured
\label{conj:thqg-shadow-depth-code-distance}
\index{shadow depth!filtered recovery models}
\textit{Type signature.} Quadrant: Open. Presentation: chain-level
(formal MC tower) plus external Hilbert error model. Level: $5$
(scalar shadow tower) lifted to physical recovery. Hypothesis
package: chirally Koszul $\cA$ in the standard-landscape scope;
fixed physical error model whose erasures match the non-scalar
shadow degrees.
\smallskip\noindent
Let a physical error model be supplied whose erasures are exactly the
non-scalar shadow degrees of the MC tower.  Then the formal number of
filtered recovery slots is
\[
\#\{\,r\ge 3:\Sh_r(\Theta_\cA)\neq 0\,\}.
\]
On a pure shadow-depth archetype this number is $r_{\max}(\cA)-2$
for finite $r_{\max}$ and countably infinite for class~$\mathbf M$.
The physical inner product, error algebra, and recovery maps determine
the associated Hilbert-space code distance.
\end{conjecture}

\begin{theorem}[Shadow depth controls the formal shadow slots;
\ClaimStatusConditional]
\label{thm:hc-shadow-redundancy}
\index{shadow depth!formal slots}
\textit{Type signature.} Quadrant: Open. Presentation: chain-level
(MC obstruction tower). Level: $5$ (scalar shadow algebra
$\cA^{\mathrm{sh}}_{r,\bullet}$) sourced from the level-$2$ bar
intrinsic MC element $D_\cA - \dzero$. Hypothesis package: chirally
Koszul $\cA$; standard-landscape scope of
Theorem~\textup{\ref{thm:shadow-archetype-classification}};
recursive existence theorem
(Theorem~\textup{\ref{thm:recursive-existence}}); shadow growth bound
(Theorem~\textup{\ref{thm:shadow-radius}}) for class~M convergence.
\smallskip\noindent
For a chirally Koszul algebra $\cA$ in the computed standard
landscape scope of
Theorem~\textup{\ref{thm:shadow-archetype-classification}}, the shadow
depth $r_{\max}(\cA)$ controls the non-scalar slots in the formal MC
obstruction tower as follows.

\begin{enumerate}[label=\textup{(\roman*)}]
\item \emph{Fundamental datum.}
 The scalar-level shadow $\kappa(\cA)$
 \textup{(}degree~$2$\textup{)} is the first modular datum in the
 tower: it is the genus-$1$ scalar obstruction of Theorem~D.

\item \emph{Recursive lifting.}
 For each degree $r \ge 3$, the shadow
 $\Sh_r(\Theta_\cA) \in \cA^{\mathrm{sh}}_{r,\bullet}$
 \textup{(}as a cohomology class in the shadow algebra\textup{)}
 is the degree-$r$ projection of the canonical bar-intrinsic MC
 element $D_\cA-\dzero$.
 The recursive existence theorem
 \textup{(}Theorem~\textup{\ref{thm:recursive-existence}}\textup{)}
 gives a lift from degree $<r$ to degree~$r$ after the obstruction
 class in $H^2$ vanishes; choices of representative are controlled by
 the degree-$r$ $H^1$.  A contraction, gauge convention, or physical
 recovery map turns a degree-$r$ shadow into a formal recovery slot.

\item \emph{Slot count.}
 The number of non-scalar shadow slots on the pure archetype is:
 \begin{center}
 \renewcommand{\arraystretch}{1.3}
 \begin{tabular}{ccll}
 \textbf{Class} & $r_{\max}$ & \textbf{Non-scalar slots}
 & \textbf{Prototype} \\
 \midrule
 G & $2$ & $0$ & Heisenberg, lattice \\
 L & $3$ & $1$ & Affine Kac--Moody \\
 C & $4$ & $2$ & $\beta\gamma$ \\
 M & $\infty$ & $\aleph_0$ \textup{(}countably infinite\textup{)}
 & Virasoro, $W_N$
 \end{tabular}
 \end{center}
 For class~M with shadow radius $\rho(\cA) < 1$, the scalar shadow
 series converges absolutely.
\end{enumerate}
These data determine formal slot count.  The universal degree statement
is that the scalar tower begins at degree~$2$; the physical error
package determines Hilbert-space code distance.
\end{theorem}

\begin{proof}
\emph{Part (i).} The scalar shadow $\kappa(\cA)$ is the
genus-$1$ obstruction class (Theorem~D), determining the
modular curvature on the uniform-weight scalar lane.  Degree~$2$ is
the initial degree of this shadow tower.

\emph{Part (ii).} The recursive existence theorem
(Theorem~\ref{thm:recursive-existence}) constructs the
shadow obstruction tower $\Theta_\cA = \varprojlim \Theta_\cA^{\le r}$
by solving the MC equation at each degree. At degree $r \ge 3$,
the equation
$d(\Theta^{(r)}) + \tfrac{1}{2}\sum_{r_1+r_2=r}
[\Theta^{(r_1)}, \Theta^{(r_2)}] = 0$
determines the \emph{shadow class}
$[\Theta^{(r)}] \in H^0(\operatorname{gr}^r \fg^{\mathrm{mod}})$
from $\{\Theta^{(s)}\}_{s < r}$.
The representatives $\Theta^{(r)}$ form the gauge torsor controlled
by~$H^1$; $H^2$ carries the obstruction.
The bar-intrinsic representative $D_\cA-\dzero$ supplies the canonical
choice used in the monograph.  A physical error package promotes this
formal lifting slot to a Hilbert-space recovery datum.

\emph{Part (iii).}
Class~G: $\Sh_r = 0$ for $r \ge 3$; zero non-scalar slots.
Class~L: $\Sh_3 \ne 0$, $\Sh_r = 0$ for $r \ge 4$; one non-scalar slot.
Class~C: $\Sh_4 \ne 0$ (and possibly $\Sh_3$), $\Sh_r = 0$
for $r \ge 5$; two non-scalar slots.
Class~M: $\Sh_r \ne 0$ for all $r$; countably many non-scalar slots.
Convergence for class~M follows from the shadow growth rate
bound $|\Sh_r| = O(\rho(\cA)^r \cdot r^{-5/2})$
(Theorem~\ref{thm:shadow-radius}), with absolute convergence
for $\rho < 1$.
\end{proof}

% ======================================================================
%
% 4. THE 12-FOLD DICTIONARY
%
% ======================================================================

\section{The Koszulness--code dictionary}
\label{sec:hc-dictionary}
\index{Koszulness--code dictionary|textbf}

\begin{proposition}[Koszulness--code dictionary;
\ClaimStatusConditional]
\label{prop:hc-dictionary}
\index{Koszulness!code-theoretic meaning}
\textit{Type signature.} Quadrant: Open. Presentation: dictionary
(structural readings of the twelve entries in
Theorem~\ref{thm:koszul-equivalences-meta}). Levels:
$1\!\leftrightarrow\!2$ for K1--K5 and K9--K10; categorical
resolution for K6; factorization homology for K7; level~$3$ for K8;
level~$5$ for K11; geometric purity for K12. Hypothesis package:
$H_{\mathrm{CL}}$ for the intrinsic core; the displayed monadic,
factorization-homology, $H_H$, shifted-symplectic, and mixed-Hodge
packages for K6, K7, K8, K11, and K12, respectively.
\smallskip\noindent
The twelve entries consist of seven intrinsic diagnostics for
$q_\cA$, one monadic consequence, one conditional
factorization-homology criterion, one $H_H$-conditional Hochschild
consequence, one shifted-symplectic criterion, and one mixed-Hodge
purity implication.  The third column records the algebraic datum
that enters a code model; Hilbert-space QEC arises from the additional
physical package of Definition~\ref{def:hc-typed-recovery-data}.

\begin{center}
\renewcommand{\arraystretch}{1.4}
\small
\begin{tabular}{l@{\;}l@{\;\;}l@{\;\;}c}
& \textbf{Koszulness condition}
& \textbf{Code property}
& \textbf{Status} \\
\midrule
K1 & Koszul twisting morphism
 & presentation-level encoding datum
 & intrinsic \\
K2 & PBW diagonal collapse
 & graded quadratic compression
 & intrinsic \\
K3 & off-diagonal $A_\infty$ formality
 & minimal compressed model
 & intrinsic \\
K4 & Ext diagonal concentration
 & diagonal defect groups
 & intrinsic \\
K5 & $q_\cA$ and $p_\cA$ quasi-isomorphisms \textup{(}Thm~B\textup{)}
 & \textbf{exact quadratic recovery}
 & exact \\
K6 & Barr--Beck--Lurie comparison
 & categorical recovery
 & cond. \\
K7 & factorization-homology concentration
 & global homological signature
 & cond. \\
K8 & $H_H(\cA;S)$-conditional $\ChirHoch$ support
 & family-indexed closed-sector support
 & $\Rightarrow$ \\
K9 & Kac--Shapovalov nonvanishing
 & non-degenerate Gram matrix
 & intrinsic \\
K10 & FM boundary acyclicity
 & collision-local recovery
 & intrinsic \\
K11 & Lagrangian criterion
 & half-dimensional Lagrangian subspace
 & cond. \\
K12 & $\mathcal D$-module purity
 & geometric sufficient condition for K10
 & $\Rightarrow$
\end{tabular}
\end{center}
Here ``intrinsic'' means equivalent under $H_{\mathrm{CL}}$ and its
PBW detection package; ``cond.'' imports the package displayed in the
corresponding row; $\Rightarrow$ records a one-way implication.
\end{proposition}

\begin{proof}
The ordering is the ordering of
Theorem~\ref{thm:koszul-equivalences-meta}.  Its intrinsic circuit is
K1--K5 together with K9--K10.  The fundamental twisting theorem
\cite[Theorem~2.3.2]{LodayVallette12} and the quadratic criterion
\cite[Theorem~3.4.6]{LodayVallette12} identify K5 with
\[
q_\cA\colon\cA^{\mathrm i}\xrightarrow{\sim}B_X(\cA),
\qquad
p_\cA\colon\Omega_X(\cA^{\mathrm i})\xrightarrow{\sim}\cA.
\]
Thus K5 is exact recovery from the quadratic compression.  K6 is the
monadic consequence under conservativity and totalization.  K7 and K8
are the factorization-homology and derived-centre consequences under
their named comparison packages.  K9 says that the relevant Gram form
is non-degenerate; a physical inner product transports this datum to a
Hilbert-space Gram matrix.  K10 is the collision-stratum form of the
intrinsic acyclicity test.

\emph{K11 (Lagrangian criterion).}
Under perfectness and non-degeneracy, Theorem~C gives
$\mathbf C_g(\cA)=\mathbf Q_g(\cA)\oplus\mathbf Q_g(\cA^!)$ with
both summands Lagrangian; in finite dimension this forces
$\dim\mathbf Q_g(\cA)=\frac12\dim\mathbf C_g(\cA)$.  This is the
Verdier-symplectic rate statement.  K12 supplies a geometric sufficient
condition for K10.  Proposition~\ref{prop:hc-theorem-h-scope} and
Definition~\ref{def:hc-typed-recovery-data} record the comparison maps
that carry K8 and K11 toward physical observables and Hilbert-space
recovery.
\end{proof}

\begin{remark}[Algebraic parameters supplied by the model]
\label{rem:hc-universal-parameters}
\index{code rate!universal $1/2$}
\index{code distance!physical input}
The algebraic model supplies three parameters, each with its own
scope:
\begin{enumerate}[label=\textup{(\roman*)}]
\item \emph{Verdier rate.}
 Under the finite-dimensional Lagrangian hypotheses of Theorem~C,
 $\dim \cC=\tfrac12\dim\cH$.
\item \emph{First degree.}
 The shadow tower begins at degree~$2$ with the scalar
 $\kappa(\cA)$.
\item \emph{Formal shadow slots.}
 On a pure archetype, the number of non-scalar slots is
 $r_{\max}(\cA)-2$ for finite $r_{\max}$ and countably infinite for
 class~$\mathbf M$.
\end{enumerate}
The structure is symplectic in the Verdier sense.  A physical
Hilbert-space realization and an error algebra determine the
stabilizer group and operational code distance.
\end{remark}

% ======================================================================
%
% 5. STANDARD LANDSCAPE CODE CENSUS
%
% ======================================================================

\section{Code census for the standard landscape}
\label{sec:hc-census}
\index{standard landscape!code census|textbf}

\begin{theorem}[Standard landscape code census;
\ClaimStatusConditional]
\label{thm:hc-census}
\textit{Type signature.} Quadrant: Open. Presentation: chain-level.
Level: $1\!\leftrightarrow\!2$ (quadratic recognition) plus level-$5$
shadow-slot count. Hypothesis package: a Theorem~B recognition datum
$\bigl(\cA^{\mathrm i},q_\cA,H_{\mathrm{CL}}\bigr)$ for each row;
shadow archetype classification
(Theorem~\textup{\ref{thm:shadow-archetype-classification}});
shadow-slot tabulation
(Theorem~\textup{\ref{thm:hc-shadow-redundancy}}); for class M,
absolute convergence requires
$\rho(\cA) < 1$
(Theorem~\textup{\ref{thm:shadow-radius}}).
\smallskip\noindent
Each row carrying this recognition datum admits exact quadratic
recovery
\[
\cA^{\mathrm i}\xrightarrow{\sim}B_X(\cA),
\qquad
\Omega_X(\cA^{\mathrm i})\xrightarrow{\sim}\cA.
\]
The universal counit belongs to Theorem~A for every augmented row.
The formal shadow slots and heuristic physical labels are:

\begin{center}
\renewcommand{\arraystretch}{1.3}
\small
\begin{tabular}{l@{\;\;}c@{\;\;}c@{\;\;}c@{\;\;}l}
\textbf{Family} & \textbf{Class} & \textbf{Slots}
& \textbf{Algebraic type} & \textbf{Heuristic physical label} \\
\midrule
Heisenberg $\mathcal{H}_\kappa$
 & G & $0$ & scalar & abelian CS \\
Lattice $V_\Lambda$
 & G & $0$ & scalar & lattice gauge \\
Affine $\widehat{\mathfrak{g}}_k$
 & L & $1$ & cubic & non-abelian CS \\
$\beta\gamma$
 & C & $2$ & quartic/contact & topological B-model \\
Virasoro $\mathrm{Vir}_c$
 & M & $\aleph_0$ & infinite shadow tower & 3d gravity \\
$W_N$
 & M & $\aleph_0$ & infinite shadow tower & higher-spin gravity
\end{tabular}
\end{center}
For class~M families, the scalar shadow series converges absolutely
when $\rho(\cA) < 1$.
\end{theorem}

\begin{proof}
The recognition datum supplies Theorem~\ref{thm:bar-cobar-inversion-qi}
for each asserted row.  The shadow slots follow from
the shadow archetype classification
(Theorem~\ref{thm:shadow-archetype-classification}) and
Theorem~\ref{thm:hc-shadow-redundancy}.
\end{proof}

\begin{remark}[Heuristic gravitational reading of the shadow slots]
\label{rem:hc-gravitational-code-complexity}
\index{gravitational complexity!code complexity}
The following physical reading is heuristic.  Its promotion package
consists of the derived-centre/BRST comparison and Hilbert-space
recovery maps:
\begin{itemize}
\item \emph{Scalar type} (G):
 The shadow tower has zero non-scalar slots.  In the Gaussian
 gravitational reading, the scalar entropy equals its scalar term
 (Chapter~\ref{chap:entanglement-modular-koszul},
 Proposition~\ref{prop:ent-complexity-classification}).

\item \emph{Cubic type} (L):
 One non-scalar slot, the cubic shadow, is available.  In the
 tree-level reading this is the algebraic place where a three-point
 interaction can be compared with lower-degree data.

\item \emph{Contact type} (C):
 Two non-scalar slots, cubic and quartic, are available.
 The quartic contact invariant
 $\mathfrak{Q}^{\mathrm{contact}}$
 is the second formal shadow datum.

\item \emph{Infinite shadow type} (M):
 There is one formal slot at each degree.  The MC equation supplies
 the recursive algebraic constraint.  Convergence for $\rho<1$
 establishes convergence of the scalar shadow series; the physical
 comparison and recovery maps furnish error-correction content.
\end{itemize}
\end{remark}
